% ============================================================
% CircMAC: Circular-aware Multi-branch Architecture
%          for circRNA-miRNA Interaction Prediction
% Paper Draft — Introduction / Method / Experiments
% ============================================================

\documentclass[10pt,twocolumn]{article}
\usepackage[margin=1in]{geometry}
\usepackage{amsmath,amssymb,bm}
\usepackage{graphicx}
\usepackage{booktabs}
\usepackage{multirow}
\usepackage{xcolor}
\usepackage{hyperref}
\usepackage{enumitem}
\usepackage{algorithm}
\usepackage{algpseudocode}

\newcommand{\todo}[1]{\textcolor{red}{[TODO: #1]}}
\newcommand{\circmac}{\textsc{CircMAC}}
\newcommand{\dcirc}{d_{\text{circ}}}

\title{\textbf{CircMAC: Circular-aware Multi-branch Architecture\\
       for circRNA--miRNA Interaction Prediction}}

\author{
  \todo{Author 1} \quad \todo{Author 2} \\
  \todo{Affiliation} \\
  \texttt{\todo{email}}
}

\date{}

\begin{document}
\maketitle

% ============================================================
\begin{abstract}
Circular RNAs (circRNAs) are covalently closed non-coding RNA molecules
that regulate gene expression by sequestering microRNAs (miRNAs) through
a sponge mechanism. Accurate prediction of circRNA--miRNA binding sites
is critical for understanding gene regulatory networks and disease
mechanisms. However, all existing deep learning methods model circRNAs as
linear sequences, ignoring the Back-Splice Junction (BSJ) where the 3'
and 5' ends covalently connect---thereby fragmenting binding patterns that
span this biologically critical region.
%
We propose \circmac{} (\textbf{Circ}ular-aware \textbf{M}ulti-branch
\textbf{A}ttention and \textbf{C}onvolutional architecture), a lightweight
($\approx$3M parameters) encoder designed specifically for the circular
topology of circRNA. \circmac{} integrates three complementary branches:
(i)~Transformer attention with a \emph{circular relative position bias}
($\dcirc(i,j)=\min(|i{-}j|,\,L{-}|i{-}j|)$),
(ii)~Mamba for efficient long-range sequential modeling, and
(iii)~depthwise CNN for local nucleotide pattern extraction---all fused
by a learned adaptive router.
%
For self-supervised pretraining on unlabeled circRNA data, we introduce
\emph{Circular Permutation Contrastive Learning} (CPCL), which uses
circular rotations of a sequence as positive pairs to learn
rotation-invariant representations, combined with standard
Masked Language Modeling (MLM).
At inference, circRNA representations interact with miRNA via
cross-attention, and per-nucleotide binding sites are decoded by a
multi-scale convolutional head.
%
Experiments on \todo{dataset stats} samples show that \circmac{}
achieves competitive or superior performance compared to large pretrained
RNA language models (\textgreater95M parameters) while using $\sim$30$\times$
fewer parameters and covering 100\% of sequence lengths.
\end{abstract}

% ============================================================
\section{Introduction}
\label{sec:intro}

Circular RNAs (circRNAs) have emerged as key regulators of gene expression,
acting as competitive endogenous RNAs (ceRNAs) that sequester microRNAs
through a molecular sponge mechanism~\cite{hansen2013natural}.
Unlike linear RNAs, circRNAs are formed by back-splicing, where a
downstream splice donor is joined to an upstream splice acceptor,
creating a covalently closed loop at the \emph{Back-Splice Junction} (BSJ).
This circular topology confers remarkable stability and unique functional
properties~\cite{jeck2013circular}.
Dysregulation of circRNA--miRNA interactions has been implicated in cancer,
neurological disorders, and cardiovascular disease~\todo{cite},
making their accurate computational prediction a pressing challenge.

\paragraph{Limitations of existing methods.}
Prior computational approaches to circRNA--miRNA interaction prediction
range from thermodynamic models (RNAhybrid~\cite{rehmsmeier2004fast},
miRanda~\cite{enright2003microrna}) to deep learning methods using
CNNs~\todo{cite}, RNNs~\todo{cite}, and Transformers~\todo{cite}.
Despite this progress, \emph{all} existing methods commit a fundamental
modeling error: they treat the circRNA sequence as linear.
This linearization produces three concrete failure modes:
%
\begin{enumerate}[leftmargin=*,itemsep=2pt]
  \item \textbf{Artificial boundary at BSJ.}
    The junction where positions $1$ and $L$ are physically adjacent is
    treated as two distant sequence ends, breaking binding motifs that
    span this region.
  \item \textbf{Non-invariant positional encodings.}
    Standard absolute or relative position encodings assign maximally
    distant representations to BSJ-adjacent positions, when they should
    be closest.
  \item \textbf{Fragmented receptive field.}
    Convolutional filters and attention windows cannot model continuity
    across the BSJ boundary.
\end{enumerate}

\paragraph{Our approach.}
We address these limitations with \circmac{}, a architecture that
explicitly models the circular nature of circRNA at every level:
attention, convolution, and pretraining.
Additionally, we observe that large pretrained RNA language models
(RNA-FM~\cite{chen2022interpretable}: 95M params, RNA-MSM~\cite{zhang2024multiple}: 96M params)
were designed for \emph{linear} RNA and saturate sequence coverage
at 1,024 nucleotides, the same limit as our lightweight model.
Our key insight is that a purpose-built circular architecture can
match or outperform these large models at a fraction of the cost.

\paragraph{Contributions.} We summarize our contributions as follows:
\begin{enumerate}[leftmargin=*,itemsep=2pt]
  \item \textbf{\circmac{} Architecture:} A multi-branch encoder
    (Attention + Mamba + CNN) with circular relative position bias
    and an adaptive routing mechanism, operating in a multi-scale
    framework for computational efficiency.
  \item \textbf{Circular-aware Pretraining:}
    \emph{Circular Permutation Contrastive Learning} (CPCL), a novel
    self-supervised objective that exploits the rotational symmetry of
    circRNAs to learn BSJ-invariant representations.
  \item \textbf{Comprehensive Evaluation:}
    Six experiments covering encoder architecture, pretraining strategy,
    comparison with pretrained RNA models, ablation, interaction mechanism,
    and site head design---demonstrating \todo{X\%} improvement on binding
    site F1 over strong baselines.
\end{enumerate}

% ============================================================
\section{Related Work}
\label{sec:related}

\paragraph{circRNA--miRNA interaction prediction.}
Classical tools such as miRanda~\cite{enright2003microrna} and
RNAhybrid~\cite{rehmsmeier2004fast} rely on sequence complementarity
and thermodynamic stability, without learning from data.
Deep learning methods have since improved accuracy using
CNNs~\todo{cite}, bidirectional LSTMs~\todo{cite}, and attention
mechanisms~\todo{cite}, but invariably represent circRNA as a linear
sequence.

\paragraph{RNA foundation models.}
RNABERT~\cite{akiyama2022informative}, RNAErnie~\cite{wang2023rnasolo},
RNA-FM~\cite{chen2022interpretable}, and RNA-MSM~\cite{zhang2024multiple}
adopt the BERT-style pretraining paradigm on large corpora of linear RNA.
While powerful, these models are limited by positional encodings designed
for linear sequences and carry \textgreater85M parameters, making
fine-tuning expensive.
None were designed with circRNA-specific inductive biases.

\paragraph{State space models for sequences.}
Mamba~\cite{gu2023mamba} and its predecessor S4~\cite{gu2021efficiently}
achieve linear-time sequence modeling through selective state spaces,
offering an efficient alternative to quadratic-complexity attention.
Hymba~\cite{dong2024hymba} combines attention and SSM heads.
We build on these models but adapt them for circular sequences.

% ============================================================
\section{Method}
\label{sec:method}

\subsection{Problem Formulation}

Let $\mathbf{c} = (c_1, c_2, \ldots, c_L)$ denote a circRNA sequence of
length $L$, where the topology implies $c_L$ is physically adjacent to
$c_1$ at the BSJ.
Let $\mathbf{m} = (m_1, m_2, \ldots, m_M)$ denote a miRNA sequence.
The binding site prediction task requires learning a function:
\begin{equation}
  g(\mathbf{c}, \mathbf{m}) \;\to\; \hat{\mathbf{y}} = (\hat{y}_1, \ldots, \hat{y}_L),
  \quad \hat{y}_i \in \{0,1\}
\end{equation}
where $\hat{y}_i=1$ indicates nucleotide $c_i$ participates in a binding
site. We additionally consider the auxiliary binding prediction task
$f(\mathbf{c},\mathbf{m}) \to \{0,1\}$ (does any interaction occur?).

\subsection{Overall Framework}

The \circmac{} framework consists of three components:
(1)~a circular-aware \circmac{} encoder for the circRNA,
(2)~a frozen RNABERT encoder for the miRNA,
and (3)~a cross-attention interaction module followed by a convolutional
site prediction head (Figure~\ref{fig:overview}).
The circRNA encoder is the focus of this work; miRNA is encoded by a
fixed pretrained model as it is short ($M \leq 25$) and linear.

\subsection{CircMAC Encoder}

\paragraph{Multi-scale structure.}
Each \circmac{} block operates at half-resolution to reduce computational
cost and increase the effective receptive field.
The encoder first downsamples the token embeddings
$\mathbf{H}^{(0)} \in \mathbb{R}^{B \times L \times D}$
via a stride-2 depthwise convolution to
$\tilde{\mathbf{H}} \in \mathbb{R}^{B \times L/2 \times D}$,
processes them through $N$ \circmac{} blocks, then upsamples back to
full resolution via linear interpolation and adds a skip connection:
\begin{equation}
  \mathbf{E}_{\text{circ}} = \text{Upsample}\!\left(\tilde{\mathbf{H}}^{(N)}\right)
    + \mathbf{H}^{(0)}.
\end{equation}

\paragraph{CircMAC block.}
Each block receives a hidden state $\mathbf{x} \in \mathbb{R}^{B \times L/2 \times D}$,
normalizes it with RMSNorm, projects it to queries/keys/values and a
``base'' feature via a shared linear layer, then passes the base
through three parallel branches:

\subsubsection{Branch 1: Circular-aware Attention}

Standard relative position bias penalizes token pairs by
$d(i,j) = |i-j|$, assigning positions $1$ and $L$ the maximum distance.
We replace this with the \emph{circular distance}:
\begin{equation}
  \dcirc(i,j) = \min\!\left(|i-j|,\; L - |i-j|\right),
  \label{eq:circ_bias}
\end{equation}
so that BSJ-adjacent positions $(i{=}1, j{=}L)$ receive the same small
distance as consecutive positions.
The attention score matrix is computed as:
\begin{equation}
  A_{ij} = \frac{\mathbf{q}_i \cdot \mathbf{k}_j}{\sqrt{d}} - \lambda \cdot \dcirc(i,j),
\end{equation}
where $\lambda$ is a learnable per-head slope parameter.
This \emph{circular relative bias} enables attention across the BSJ
without any change to token embeddings.

\subsubsection{Branch 2: Mamba (Sequential)}

The base feature is passed through a Mamba selective SSM block
\cite{gu2023mamba} with state dimension $d_{\text{state}}{=}16$ and
expansion factor $\text{expand}{=}2$ (inner dimension 256).
Mamba captures long-range sequential dependencies in $O(L)$ time.

\subsubsection{Branch 3: Depthwise Circular CNN}

To model local nucleotide patterns, we apply a depthwise convolution
with kernel size $k{=}7$. The key modification is \emph{circular padding}:
instead of zero-padding, we wrap the sequence ends:
\begin{equation}
  \tilde{\mathbf{x}} = [x_{L-p+1}, \ldots, x_L,\; x_1, \ldots, x_L,\; x_1, \ldots, x_p],
\end{equation}
where $p = (k-1)/2 = 3$. This allows the convolutional kernel to see
continuity across the BSJ, capturing motifs that span it.

\subsubsection{Adaptive Router}

Rather than using fixed branch weights, we employ a lightweight
sequence-level router. The three branch outputs are concatenated,
mean-pooled over the length dimension, and passed through an MLP:
\begin{align}
  \mathbf{g} &= \text{Softmax}\!\left(\,
    W_2\,\text{GELU}(W_1\,\overline{\mathbf{x}}_{\text{cat}} + b_1) + b_2
  \right) \in \mathbb{R}^3, \\
  \mathbf{y} &= g_1\,\mathbf{y}_{\text{attn}}
             + g_2\,\mathbf{y}_{\text{mamba}}
             + g_3\,\mathbf{y}_{\text{cnn}},
\end{align}
where $W_1 \in \mathbb{R}^{3D \times D}$ and $W_2 \in \mathbb{R}^{D \times 3}$.
The router adapts the contribution of each branch to the characteristics
of each input sequence.

\subsection{circRNA--miRNA Cross-Attention Interaction}
\label{sec:interaction}

After encoding, the circRNA representation
$\mathbf{E}_{\text{circ}} \in \mathbb{R}^{B \times L \times D}$ and
miRNA representation $\mathbf{E}_{\text{mirna}} \in \mathbb{R}^{B \times M \times D}$
are fused via cross-attention, where the circRNA serves as queries:
\begin{equation}
  \text{Context} = \text{softmax}\!\left(
    \frac{(\mathbf{E}_{\text{circ}} W_Q)(\mathbf{E}_{\text{mirna}} W_K)^\top}{\sqrt{d_k}}
  \right) (\mathbf{E}_{\text{mirna}} W_V).
\end{equation}
The context is concatenated with $\mathbf{E}_{\text{circ}}$ and
projected back to dimension $D$:
\begin{equation}
  \mathbf{E}_{\text{fused}}
    = W_O \,[\mathbf{E}_{\text{circ}}\,\|\,\text{Context}]
    \in \mathbb{R}^{B \times L \times D}.
\end{equation}
This design allows every circRNA nucleotide to directly attend to the
full miRNA sequence, enabling site-level interaction modeling.

\subsection{Site Prediction Head}

The fused representation passes through a multi-kernel CNN to capture
binding patterns at multiple scales:
\begin{equation}
  \mathbf{F} = \frac{1}{3}\sum_{k \in \{3,5,7\}}
    \text{Conv1d}_k(\mathbf{E}_{\text{fused}}) + \mathbf{E}_{\text{fused}},
\end{equation}
followed by two pointwise convolutions
$D \to D/2 \to 2$ with BatchNorm and GELU activation, producing
per-position logits $\hat{\mathbf{y}} \in \mathbb{R}^{B \times L \times 2}$.

\subsection{Self-supervised Pretraining}
\label{sec:pretrain}

We pretrain \circmac{} on a corpus of \todo{N} unlabeled circRNA
sequences using multiple self-supervised objectives, combined with
learned uncertainty weighting~\cite{kendall2018multi}:
\begin{equation}
  \mathcal{L} = \sum_k \frac{1}{\sigma_k^2} \mathcal{L}_k + \log \sigma_k^2,
\end{equation}
where $\sigma_k$ are learnable per-task parameters.

\paragraph{MLM (Masked Language Modeling).}
We randomly mask 15\% of nucleotide tokens and train the model to
reconstruct them, following \citet{devlin2019bert}.

\paragraph{CPCL (Circular Permutation Contrastive Learning).}
circRNAs are rotationally symmetric: the sequence starting at position
$k$ represents the same molecule as the canonical sequence starting at
position 1. We exploit this as a free source of positive pairs.
For a sequence $\mathbf{c}$ and a randomly shifted version
$\mathbf{c}' = (c_{k+1}, \ldots, c_L, c_1, \ldots, c_k)$,
we maximize their representation agreement using NT-Xent loss:
\begin{equation}
  \mathcal{L}_{\text{CPCL}} = -\log
    \frac{\exp(\text{sim}(\mathbf{z}, \mathbf{z}') / \tau)}
         {\sum_{j \neq i} \exp(\text{sim}(\mathbf{z}, \mathbf{z}_j) / \tau)},
\end{equation}
where $\mathbf{z}, \mathbf{z}'$ are mean-pooled representations and
$\tau=0.1$ is the temperature.
CPCL encourages \circmac{} to learn representations invariant to the
arbitrary choice of BSJ linearization start point.

\paragraph{Additional objectives.}
We also experiment with Next Token Prediction (NTP), Secondary Structure
Prediction (SSP), and base-pair matrix reconstruction using a circular
pairing head; the best combination is determined by Experiment 2
(Section~\ref{sec:exp2}).

% ============================================================
\section{Experiments}
\label{sec:experiments}

\subsection{Experimental Setup}
\label{sec:setup}

\paragraph{Dataset.}
We use a circRNA--miRNA interaction dataset compiled from \todo{source},
consisting of \todo{N} sequence pairs with token-level binding site
annotations. The dataset is split into training (\todo{N$_{\text{train}}$}
samples), validation (\todo{N$_{\text{val}}$}), and test
(\todo{N$_{\text{test}}$}).
For pretraining, we use the \texttt{df\_circ\_ss\_5} corpus containing
\todo{$\sim$149K} unlabeled circRNA sequences with secondary structure
annotations (up to $L{=}1022$ nucleotides).

\paragraph{Evaluation metrics.}
For the \emph{binding} task (binary): Accuracy, F1, Precision, Recall,
MCC, AUROC, AUPRC.
For the \emph{site} task (token-level): Accuracy, F1, Precision, Recall.
All results are averaged over 3 random seeds (seeds 1, 2, 3) and
reported as mean$\pm$std.

\paragraph{Implementation details.}
\circmac{} uses $D{=}128$, $N{=}6$ layers, $H{=}8$ attention heads.
Fine-tuning: AdamW, lr $= 10^{-4}$, batch size 32 (all experiments),
150 epochs with early stopping (patience 20).
Pretraining: AdamW, lr $= 5{\times}10^{-4}$, batch size 64,
200 epochs with early stopping (patience 30).
We encode miRNA with frozen RNABERT~\cite{akiyama2022informative}.
All experiments use PyTorch on NVIDIA GPUs.

\paragraph{Baselines.}
We compare against: (1)~sequence encoders trained from scratch
(LSTM, Transformer, Mamba, Hymba); (2)~general-purpose pretrained
RNA models (RNABERT, RNAErnie, RNA-FM, RNA-MSM) in both frozen and
trainable settings; and (3)~\circmac{} without pretraining.

\subsection{Exp 1: Encoder Architecture Comparison}
\label{sec:exp1}

We compare five encoder architectures trained from scratch under
identical conditions (Table~\ref{tab:exp1}).
All models use $D{=}128$, 6 layers, cross-attention interaction,
and the same training configuration.

\begin{table}[h]
\centering
\caption{Encoder architecture comparison on the \emph{sites} task.
All models trained from scratch (mean$\pm$std, 3 seeds).}
\label{tab:exp1}
\resizebox{\columnwidth}{!}{%
\begin{tabular}{lcccc}
\toprule
\textbf{Model} & \textbf{Type} & \textbf{Acc} & \textbf{F1} & \textbf{Recall} \\
\midrule
LSTM           & RNN           & \todo{--}    & \todo{--}   & \todo{--}       \\
Transformer    & Attention     & \todo{--}    & \todo{--}   & \todo{--}       \\
Mamba          & SSM           & \todo{--}    & \todo{--}   & \todo{--}       \\
Hymba          & Hybrid        & \todo{--}    & \todo{--}   & \todo{--}       \\
\midrule
\circmac{}     & \textbf{Ours} & \todo{--}    & \todo{\textbf{--}} & \todo{\textbf{--}} \\
\bottomrule
\end{tabular}%
}
\end{table}

\subsection{Exp 2: Pretraining Strategy}
\label{sec:exp2}

We ablate over seven pretraining configurations
(Table~\ref{tab:exp2}) to identify which combination of objectives
best transfers to the binding site task.

\begin{table}[h]
\centering
\caption{Effect of pretraining strategy on \emph{sites} F1.
\circmac{} fine-tuned after each pretraining configuration.}
\label{tab:exp2}
\resizebox{\columnwidth}{!}{%
\begin{tabular}{lcc}
\toprule
\textbf{Pretraining Config} & \textbf{Objectives} & \textbf{Sites F1} \\
\midrule
No pretrain      & ---                       & \todo{--} \\
MLM              & MLM                       & \todo{--} \\
MLM + NTP        & MLM, NTP                  & \todo{--} \\
MLM + SSP        & MLM, SSP                  & \todo{--} \\
MLM + CPCL       & MLM, \textbf{CPCL}        & \todo{--} \\
MLM + Pairing    & MLM, Pairing              & \todo{--} \\
MLM+NTP+CPCL+Pair& MLM, NTP, CPCL, Pairing   & \todo{--} \\
\bottomrule
\end{tabular}%
}
\end{table}

\subsection{Exp 3: Comparison with Pretrained RNA Models}
\label{sec:exp3}

We compare \circmac{}-PT (pretrained with the best config from Exp~2)
against four general-purpose RNA language models in frozen and
fine-tunable settings (Table~\ref{tab:exp3}).
For fairness, each model is evaluated at its own maximum sequence length;
models limited by positional encoding are noted with $\dagger$.

\begin{table*}[ht]
\centering
\caption{Comparison with pretrained RNA models on the \emph{sites} task.
Each model tested at its maximum supported sequence length.
BS = batch size; Cov.\ = fraction of test sequences fully covered.}
\label{tab:exp3}
\begin{tabular}{llcccccc}
\toprule
\textbf{Model} & \textbf{Mode} & \textbf{Params} & \textbf{max\_len} & \textbf{Cov.}
  & \textbf{Acc} & \textbf{F1} & \textbf{Recall} \\
\midrule
RNABERT$^\dagger$  & Frozen      & 0.5M  & 440  & 37.6\% & \todo{--} & \todo{--} & \todo{--} \\
RNABERT$^\dagger$  & Trainable   & 0.5M  & 440  & 37.6\% & \todo{--} & \todo{--} & \todo{--} \\
RNAErnie$^\dagger$ & Frozen      & 86M   & 510  & 48.2\% & \todo{--} & \todo{--} & \todo{--} \\
RNAErnie$^\dagger$ & Trainable   & 86M   & 510  & 48.2\% & \todo{--} & \todo{--} & \todo{--} \\
RNA-FM             & Frozen      & 95M   & 1024 & 100\%  & \todo{--} & \todo{--} & \todo{--} \\
RNA-FM             & Trainable   & 95M   & 1024 & 100\%  & \todo{--} & \todo{--} & \todo{--} \\
RNA-MSM            & Frozen      & 96M   & 1024 & 100\%  & \todo{--} & \todo{--} & \todo{--} \\
RNA-MSM            & Trainable   & 96M   & 1024 & 100\%  & \todo{--} & \todo{--} & \todo{--} \\
\midrule
\circmac{}-PT      & Fine-tuned  & \textbf{3M} & 1022 & 100\%
  & \todo{--} & \todo{\textbf{--}} & \todo{\textbf{--}} \\
\bottomrule
\end{tabular}
\begin{flushleft}
\small $^\dagger$ Limited by fixed positional encoding; sequences longer
than max\_len are truncated.
\end{flushleft}
\end{table*}

\subsection{Exp 4: Ablation Study}
\label{sec:exp4}

We systematically remove each component of \circmac{} to quantify
its individual contribution (Table~\ref{tab:ablation}).

\begin{table}[h]
\centering
\caption{Ablation study on the \emph{sites} task. Each row removes
one component from the full \circmac{} model.}
\label{tab:ablation}
\resizebox{\columnwidth}{!}{%
\begin{tabular}{lcc}
\toprule
\textbf{Configuration} & \textbf{F1} & \textbf{$\Delta$F1} \\
\midrule
Full \circmac{}              & \todo{--} & --- \\
\midrule
w/o Attention branch         & \todo{--} & \todo{--} \\
w/o Mamba branch             & \todo{--} & \todo{--} \\
w/o CNN branch               & \todo{--} & \todo{--} \\
w/o Circular rel.\ bias      & \todo{--} & \todo{--} \\
\midrule
Attention only               & \todo{--} & \todo{--} \\
Mamba only                   & \todo{--} & \todo{--} \\
CNN only                     & \todo{--} & \todo{--} \\
\bottomrule
\end{tabular}%
}
\end{table}

\subsection{Exp 5: Interaction Mechanism}
\label{sec:exp5}

We compare three strategies for fusing circRNA and miRNA representations:
concatenation, element-wise product, and cross-attention
(Table~\ref{tab:interaction}).

\begin{table}[h]
\centering
\caption{Interaction mechanism comparison.}
\label{tab:interaction}
\begin{tabular}{lcc}
\toprule
\textbf{Interaction} & \textbf{F1} & \textbf{Recall} \\
\midrule
Concat          & \todo{--} & \todo{--} \\
Elementwise     & \todo{--} & \todo{--} \\
Cross-attention & \todo{--} & \todo{--} \\
\bottomrule
\end{tabular}
\end{table}

\subsection{Exp 6: Site Head Design}
\label{sec:exp6}

We compare the multi-kernel convolutional head against a simple
linear projection baseline (Table~\ref{tab:sitehead}).

\begin{table}[h]
\centering
\caption{Site prediction head comparison.}
\label{tab:sitehead}
\begin{tabular}{lcc}
\toprule
\textbf{Head} & \textbf{F1} & \textbf{Recall} \\
\midrule
Linear          & \todo{--} & \todo{--} \\
Conv1d (ours)   & \todo{--} & \todo{--} \\
\bottomrule
\end{tabular}
\end{table}

% ============================================================
\section{Conclusion}
\label{sec:conclusion}

We presented \circmac{}, a circular-aware architecture for
circRNA--miRNA binding site prediction that explicitly models the
circular topology of circRNA through three mechanisms: circular relative
position bias in attention, circular padding in CNN, and
Circular Permutation Contrastive Learning in pretraining.
\circmac{} achieves \todo{state-of-the-art} performance on binding site
localization while using $\sim$30$\times$ fewer parameters than large
RNA foundation models.
Our work demonstrates that incorporating domain-specific structural
inductive biases---rather than simply scaling model size---is a
productive direction for RNA sequence analysis.
Future work includes extending \circmac{} to circRNA--protein
interactions and investigating circular topology modeling for
other covalently closed biomolecules.

% ============================================================
\bibliographystyle{plain}
\begin{thebibliography}{99}

\bibitem{hansen2013natural}
T.B.\ Hansen et al., ``Natural RNA circles function as efficient
microRNA sponges,'' \textit{Nature}, 2013.

\bibitem{jeck2013circular}
W.R.\ Jeck et al., ``Circular RNAs are abundant, conserved, and
associated with ALU repeats,'' \textit{RNA}, 2013.

\bibitem{rehmsmeier2004fast}
M.\ Rehmsmeier et al., ``Fast and effective prediction of microRNA/target
duplexes,'' \textit{RNA}, 2004.

\bibitem{enright2003microrna}
A.J.\ Enright et al., ``MicroRNA targets in \textit{Drosophila},''
\textit{Genome Biology}, 2003.

\bibitem{akiyama2022informative}
M.\ Akiyama \& Y.\ Sakakibara, ``Informative RNA-base embedding for
functional RNA structural alignment and clustering,'' \textit{NAR Genomics
and Bioinformatics}, 2022.

\bibitem{wang2023rnasolo}
\todo{RNAErnie citation}

\bibitem{chen2022interpretable}
K.\ Chen et al., ``Interpretable RNA foundation model from unannotated
data for highly accurate RNA structure and function predictions,''
\textit{arXiv:2204.00300}, 2022.

\bibitem{zhang2024multiple}
\todo{RNA-MSM citation}

\bibitem{gu2023mamba}
A.\ Gu \& T.\ Dao, ``Mamba: Linear-time sequence modeling with selective
state spaces,'' \textit{arXiv:2312.00752}, 2023.

\bibitem{gu2021efficiently}
A.\ Gu et al., ``Efficiently modeling long sequences with structured
state spaces,'' \textit{ICLR}, 2022.

\bibitem{dong2024hymba}
\todo{Hymba citation}

\bibitem{devlin2019bert}
J.\ Devlin et al., ``BERT: Pre-training of deep bidirectional transformers
for language understanding,'' \textit{NAACL}, 2019.

\bibitem{kendall2018multi}
A.\ Kendall et al., ``Multi-task learning using uncertainty to weigh
losses for scene geometry and semantics,'' \textit{CVPR}, 2018.

\end{thebibliography}

\end{document}
